\section{性交时长、频率与"正常"的标准}

"我是不是太快了？""一周几次才正常？"——这是性咨询中最常见的两个问题，而它们的答案可能出乎意料：\textbf{没有统一标准，也不该有}。

\noindent\textbf{关于时长}：
\begin{itemize}
	\item 多项国际研究显示，从插入到射精的\textbf{中位时长约为 5--7 分钟}（不同研究给出的范围大致在 3--13 分钟之间）；
	\item 加上前戏，一次性活动通常持续 15--30 分钟甚至更久，个体差异极大；
	\item 许多人以为"半小时以上才算正常"，这多半是受影视或色情作品影响——那是表演，不是平均水平；
	\item 关键不在分钟数，而在\textbf{双方是否满意、是否有困扰}。若一方感到疼痛、疲惫或未被满足，再长的时长也算不上"好"。
\end{itemize}

\noindent\textbf{关于频率}：
\begin{itemize}
	\item 所谓"平均每周 1--2 次"只是统计上的中位数，\textbf{不是健康指标，更不是及格线}；
	\item 性满意度与频率并非线性关系——有人每周数次依然满足，有人每月一两次也其乐融融；
	\item 影响频率的因素众多：年龄、健康状况、压力与情绪、关系质量、激素水平、用药、育儿与工作负荷。频率的起伏本身很正常；
	\item 真正需要留意的是\textbf{变化}：若频率或欲望出现突然、持续的显著下降，并伴随痛苦或关系冲突，才值得就医或咨询。
\end{itemize}

\begin{tcolorbox}[colback=pink!5!white,colframe=red!50!black,title=贴心小叮咛]
不要拿"统计数字"或"别人的经验"来丈量自己的性生活。所谓"正常"，只有一种定义——\textbf{你们双方都感到舒适、满意、不被勉强}。这一点，只能由你们两个人共同商定。
\end{tcolorbox}
